% TOP_PROBLEM: {"problemId": "problem.zagier-conjecture-on-dedekind-zeta-values", "canonicalTitle": "Zagier Conjecture on Dedekind Zeta Values", "releaseRank": 271, "primaryDomain": "number_theory_arithmetic_geometry", "primaryDomainLabel": "Number theory & arithmetic geometry", "displayStatus": "open_with_solved_subcases", "releaseStatus": "open_with_solved_subcases"}
% !TeX program = lualatex
% ATTEMPT_STATUS: unresolved
% Research continuation by GPT_6_Astra_Ultra, 27 September 2026.
\documentclass[11pt]{article}
\usepackage[margin=25mm]{geometry}
\usepackage{fontspec}
\setmainfont{Latin Modern Roman}
\usepackage{amsmath,amssymb,mathtools}
\usepackage{unicode-math}
\setmathfont{Latin Modern Math}
\usepackage{xurl,hyperref}
\hypersetup{colorlinks=true,urlcolor=blue}
\setcounter{secnumdepth}{0}
\setlength{\parindent}{0pt}
\setlength{\parskip}{5pt}
\setlength{\emergencystretch}{3em}
\begin{document}
\section*{\texorpdfstring{Zagier Conjecture on Dedekind Zeta Values}{Zagier Conjecture on Dedekind Zeta Values}}\label{271}
\subsection{Definitions and mathematical statement}
Let $F$ have discriminant $D_F$, degree $d$, and signature $(r_1,r_2)$. Put $n=r_2$ for even $k\ge2$ and $n=r_1+r_2$ for odd $k$. The single-valued polylogarithm $\mathcal L_k$ is Zagier--Gangl's Bernoulli-number combination of $\operatorname{Li}_j$, taking imaginary part for even $k$ and real part for odd $k$. Extend it linearly to symbols. With $A_k(F)$ the \emph{allowable-symbol group} defined in their Section 2, the target is existence of $\xi_1,\ldots,\xi_n\in A_k(F)\otimes{\mathbb{Q}}$ such that
\[
 \zeta_F(k)\in{\mathbb{Q}}^\times\sqrt{|D_F|}\,
 \pi^{k(d-n)}\det\bigl(\mathcal L_k(\sigma_i\xi_j)\bigr)_{i,j=1}^n.
\]
Use one embedding from each complex-conjugate pair, and add all real embeddings when $k$ is odd. The empty determinant is one. The source's allowable-symbol conditions are essential and are not replaced by arbitrary formal symbols; no full regulator-lattice statement is added.
\par\smallskip\textit{Formulation and concept references:} \href{https://people.mpim-bonn.mpg.de/zagier/files/mpim/00-9/fulltext.pdf}{[S1]}.

\subsection{Short English statement}
Can every positive integral Dedekind zeta value be expressed, up to the specified rational and period factors, as a determinant of single-valued polylogarithms of allowable algebraic symbols?
\subsection{Research attempt: an allowable Gaussian symbol and the determinant gap}
\textit{GPT-6 Astra Ultra, 27 September 2026. Status: UNRESOLVED.}

The requirement that symbols be allowable is a real algebraic restriction. Choosing arbitrary complex numbers whose polylogarithms happen to give a desired determinant would not address this problem. I work out an exact example with a verified allowable symbol, explain the signature and normalization, and isolate the two different assertions needed for a general regulator-based proof.

\paragraph{Weight two over the Gaussian field.}
For $k=2$ the single-valued function is the Bloch--Wigner dilogarithm
\[
 D(z)=\operatorname{Im}\operatorname{Li}_2(z)
       +\log|z|\arg(1-z).
\]
The allowable group in Section 2 of [S1] is the kernel of
\[
 \beta_2:\mathbb Z[F]\longrightarrow\bigwedge^2 F^\times,
 \qquad [z]\longmapsto z\wedge(1-z),
                                                               \tag{1}
\]
with the specified zero values at $0,1$. Take $F=\mathbb Q(i)$ and $\xi=4[i]$. Multiplicativity in the exterior factor gives
\[
 \beta_2(\xi)=4(i\wedge(1-i))
             =i^4\wedge(1-i)=1\wedge(1-i)=0.
\]
Thus $\xi$ really belongs to $A_2(F)$, without assuming that $i\wedge(1-i)$ itself vanishes integrally. Passing to rational coefficients would permit dividing this allowable element by four, but the integral multiple avoids that issue entirely.

Since $|i|=1$, the logarithmic correction in $D(i)$ is zero. Absolute convergence of the dilogarithm series at $i$ gives
\[
 D(i)=\sum_{n\ge1}\frac{\sin(n\pi/2)}{n^2}
     =\sum_{m\ge0}\frac{(-1)^m}{(2m+1)^2}
     =\beta(2)>0.                                           \tag{2}
\]
Here $\beta$ is the Dirichlet beta function. Positivity follows already by grouping successive terms of the alternating series.

The Euler factors prove
\[
 \zeta_{\mathbb Q(i)}(s)=\zeta(s)L(s,\chi_{-4})\qquad(s>1).
                                                               \tag{3}
\]
For completeness, at an odd prime $\ell\equiv1\pmod4$ the Gaussian prime splits, giving $(1-\ell^{-s})^{-2}$; the right side has the same factor because $\chi_{-4}(\ell)=1$. For $\ell\equiv3\pmod4$ the factor is $(1-\ell^{-2s})^{-1}$, obtained on the right from $(1-\ell^{-s})^{-1}(1+\ell^{-s})^{-1}$. The ramified prime $2$ contributes $(1-2^{-s})^{-1}$ on both sides. These splitting statements can also be read from the factorization of $x^2+1$ over the residue fields.

At $s=2$, equations (2)--(3) and $\zeta(2)=\pi^2/6$ give
\[
 \zeta_{\mathbb Q(i)}(2)=\frac{\pi^2}{6}D(i).
\]
In the notation of the problem, $d=2$, $n=1$, and $|D_F|=4$. Choosing the embedding sending $i$ to $i$, the proposed determinant expression is
\[
 \sqrt{|D_F|}\,\pi^{2(d-n)}D(\xi)
 =2\pi^2\cdot4D(i)=8\pi^2D(i).
\]
Consequently the exact rational coefficient is
\[
 \boxed{\displaystyle
 \zeta_{\mathbb Q(i)}(2)=\frac1{48}
 \sqrt{|D_F|}\,\pi^{2(d-n)}D(\xi).}                         \tag{4}
\]
This verifies both allowability and the period factor in a genuine nonempty-determinant case. Changing to the conjugate embedding changes the determinant sign and is absorbed by $\mathbb Q^\times$, which allows either sign.

\paragraph{Why the prescribed rows are necessary.}
Complex conjugation satisfies
\[
 \mathcal L_k(\overline z)=(-1)^{k-1}\mathcal L_k(z).
                                                               \tag{5}
\]
This follows from the real coefficients and the choice of real or imaginary part in the defining single-valued combination. At even weight, values at real arguments vanish, and including both embeddings of a complex pair would add opposite rows. At odd weight those two rows agree. Thus counting one complex embedding from each pair, and adding real embeddings only at odd weight, removes these immediate forced dependencies.

This is a necessary dimension count, not a proof that the remaining rows are independent on allowable symbols. For example, in the Gaussian calculation the allowable element $[1]$ has $D(1)=0$. Replacing $\xi$ in (4) by $[1]$ would give determinant zero, although the zeta value is positive. The issue is therefore not merely producing allowable elements: one must produce enough with nondegenerate regulator values.

The empty-determinant convention also passes an exact boundary test. For $F=\mathbb Q$ and even $k=2m$, one has $d=1$, $n=0$, and $D_F=1$. The assertion reduces to $\zeta(2m)\in\mathbb Q^\times\pi^{2m}$, as in Euler's formula
\[
 \zeta(2m)=(-1)^{m+1}\frac{B_{2m}(2\pi)^{2m}}{2(2m)!}.
\]
There are no symbols to find in that case. Treating an empty determinant as zero would incorrectly destroy even this basic instance.

\paragraph{A conditional linear-algebra reduction.}
Let $R$ denote the map on allowable rational symbols given by the $n$ selected polylogarithm evaluations. Suppose there is an $n$-dimensional rational vector space $W$, a map $\phi:A_k(F)\otimes\mathbb Q\to W$, and a real regulator map $r:W\to\mathbb R^n$ such that $R=r\circ\phi$. Assume the following two additional facts:
\[
 \phi\text{ is surjective},\qquad
 \zeta_F(k)\in\mathbb Q^\times\sqrt{|D_F|}\pi^{k(d-n)}
       \det(r(w_1),\ldots,r(w_n))                           \tag{6}
\]
for a rational basis $w_1,\ldots,w_n$ of $W$, with nonzero determinant. Then choose preimages $\xi_j$ of the basis. Substitution into (6) proves exactly the formulation above. This reduction does not demand an integral basis or an equality of integral lattices; rational surjectivity is enough because the conclusion already permits a nonzero rational factor.

The basis dependence is equally precise. Replacing the $\xi_j$ by rational combinations through $C\in\operatorname{GL}_n(\mathbb Q)$ multiplies the determinant by $\det C\in\mathbb Q^\times$. Clearing all symbol denominators has the same harmless effect. But arbitrary \emph{real} combinations need not remain allowable rational symbols and would introduce arbitrary real factors. This is why scaling a column numerically to make the equality hold is not a proof.

A regulator theorem establishing the second assertion of (6) does not automatically establish the first. One still has to show that the appropriate regulator classes can be represented by the classical allowable polylogarithm symbols appearing here. Conversely, numerical nonvanishing of a candidate determinant alone does not establish its exact rational proportionality to the zeta value. These are separate gaps and should not be combined into a vague appeal to regulator rank.

There are major proved low-weight cases. In particular Goncharov--Rudenko's work proves the weight-four case by constructing the relevant polylogarithmic regulator comparison [E1]. That is a theorem in its stated weight; it is not a theorem for every $k$. The elementary computation (4) is an illustration and is not being claimed as new or as the extent of known results.

\paragraph{Numerical verification and its limits.}
The alternating expansion in (2) provides a rigorous check with a tail bound: after $N$ terms, the error is at most $(2N+1)^{-2}$. Our local check records rational intervals for $D(i)$ and verifies the factor $1/48$ symbolically. It does not try to infer the rational coefficient by an unbounded integer-relation search.

This distinction matters for larger determinants. Finitely many digits can exclude a proposed rational relation with a controlled denominator if rigorous error bounds are supplied; they cannot establish unrestricted rationality. Rational numbers with large denominators approximate every real number, and determinant cancellation can amplify relative error when columns are nearly dependent. Even a certified nonzero determinant would prove only independence of those columns, not the exact relation (6).

\paragraph{Outcome and first missing step.}
The completed work is an exact allowable-symbol calculation over $\mathbb Q(i)$, boundary and parity checks, and the sufficient rational-surjectivity reduction (6). The first unresolved step for this route at general weight and field is the construction of enough allowable classical symbols with the required regulator comparison. This attempt supplies neither that general construction nor a counterexample. The universal statement remains unresolved here.

\subsection{Sources}
\begin{itemize}
\item[S1]MPIM00-9, Classical and elliptic polylogarithms and special values of L-series.
\url{https://people.mpim-bonn.mpg.de/zagier/files/mpim/00-9/fulltext.pdf}.
\textit{Formulation source.}
\item[E1]A. B. Goncharov and D. Rudenko, \textit{Motivic correlators, cluster varieties and Zagier's conjecture on $\zeta_F(4)$}. Used only to identify the proved weight-four case and its regulator-comparison scope.
\url{https://arxiv.org/abs/1803.08585}.

\end{itemize}
\end{document}
